% intro.tex : 서론(국문판). 영문판 paper/manuscript/en/sections/intro.tex(2026-10-05 2차 수정판)를 문장 단위로 옮겼다.
% 수치와 인용 키는 영문판과 같다. 용어: MANUSCRIPT_SPEC 1.4 국문 정본, docs/GLOSSARY.md.

활동층 두께(active-layer thickness, ALT)는 영구동토 위 지표층이 계절에 따라 녹는 최대 깊이이며, 그 변화는 영구동토 위에 지은 기반 시설에 영향을 준다\cite{Karjalainen2019,Ran2022b}. ALT는 제한된 지점에서만 측정되며, 그 가운데 상당수가 35년에 걸친 Circumpolar Active Layer Monitoring(CALM) 자료에 있다\cite{Streletskiy2025}. 이 연구에서 평가한 다섯 지역의 ALT 측정값(라벨) 수는 러시아 동부 30개부터 알래스카 13,606개까지다(표~1). Stefan 해\cite{Stefan1891}는 ALT를 계수 $E$와 융해 도일(thawing degree-days, TDD) 제곱근의 곱으로 예측하며, $E$는 현지 융해 깊이 측정으로 보정한다\cite{Nelson1997}. 라벨이 있는 셀에서 $E$의 기하 평균은 레나델타 1.31\,cm\,($^{\circ}$C\,d)$^{-1/2}$, 러시아 서부 2.51, 티베트 고원 11.03이다(그림~1b). 새 영구동토 지역(대상 지역)에는 측정이 적으므로, 그곳 ALT 지도의 오차는 다른 지역(원천 자료)에 맞춘 모델에 달려 있다.

Stefan 식이나 Kudryavtsev 식의 출력은 기계학습(machine learning, ML)의 입력 특징으로 쓰였고\cite{Pilyugina2025,WangG2025}, $E$는 공변량으로 예측되었으며\cite{Ran2022b,ZhangC2024}, 물리 규칙으로 만든 라벨로 영점 장막(zero curtain) 검출기를 학습시킨 연구가 있다\cite{Gay2026}. 과정 모델 사전 학습과 잔차 모델링은 다른 환경 분야에서 쓰인다\cite{Read2019,Jia2021,Willard2022}. Gautam 등은 알래스카 CALM 68개 지점에서 무작위 70/30 분할 한 번으로 랜덤 포레스트와 Stefan 식을 비교하였다(시험 결정계수 0.24와 0.54)\cite{Gautam2025}. 대규모 ALT 지도는 거리 블록으로 검증되었다\cite{Aalto2018,Karjalainen2019,Ran2022a}. 군집된 공간 자료에서는 무작위 교차검증이 지도 정확도를 과대평가할 수 있고\cite{Roberts2017,Ploton2020,MeyerPebesma2022}, 편향 없는 지도 정확도를 얻으려면 공간 교차검증보다 설계 기반 추론을 쓰는 확률 표본이 필요하다는 주장이 있다\cite{Wadoux2021}. 위의 ALT 평가는 모두 라벨 조건 하나를 쓴다. 이웃 분야에서는 학습에서 뺀 지역이나 현지 관측이 적은 조건에서 모델을 평가하였고, 일부는 보정된 과정 모델과 비교하였다. 그 대상은 호수 수온\cite{Read2019,Jia2021,Willard2021}, 계측소 없는 유역\cite{Pool2019,Feng2023}, 지하 온도\cite{OMalley2026}, 공간 위험 예측\cite{Portes2026}이다. 우리가 아는 범위에서, 영구동토 ALT나 연평균 지중온도 연구 가운데 학습 자료에서 뺀 지역의 오차를 대상 라벨 수의 함수로, 같은 라벨로 재보정한 Stefan 식과 비교해 보고한 연구는 없다.

새 지역에는 대상 라벨이 쌓이는 동안 Stefan 식과 ML을 결합하는 규칙과 라벨을 더하는 절차가 필요하다. 이 연구는 세 질문을 다룬다. (1) 대상 라벨이 없을 때 물리 정보는 ML에 무엇을 더하며, 그 이득은 어디에서 오는가? (2) 대상 라벨이 10개에서 수천 개로 늘 때 ML은 같은 라벨로 재보정한 Stefan 식에 얼마나 더하며, 그 이득은 어느 공간 규모에서 생기는가? (3) 라벨이 쌓일 때 방법은 어떻게 고르고, 새 라벨은 어떤 절차로 배치해야 하는가?

각 대상 지역은 100\,km 완충 구역과 함께 원천 자료에서 빼고, 그 지역 0.5$^{\circ}$ 블록의 절반에서 라벨을 뽑아 나머지 절반에서 오차를 채점하였다. 라벨 수는 0, 3, 10, 40, 160, 320, 1000개와 쓸 수 있는 전량이다(그림~1c,d). 라벨이 많은 지역 안에서는 앞선 결과를 본 뒤 설계한 블록 홀드아웃 시험으로 재보정 모델에 대한 ML의 이득을 측정하였다. 대비는 블록 부트스트랩 95\% 신뢰구간(confidence interval, CI)과 $\pm$0.5\,cm 동등 한계로 판정하였고, 전이 분석은 내부 사전 등록하였다. 대상 라벨이 없을 때 Stefan 유사라벨(Stefan 예측을 학습 라벨로 쓴 것)은 섞은 유사라벨보다 주 4지역(레나델타, 캐나다, 러시아 서부, 러시아 동부)의 평균 제곱근 오차(root-mean-square error, RMSE) 평균을 낮추었다. 라벨 10개로 Stefan 계수를 재보정하면 원천 계수보다 이 값이 낮아졌다(오차가 낮아진 지역은 4곳 가운데 1곳). 알래스카 안에서 라벨 전량을 쓸 때 재보정 모델에 대한 ML 이득의 거의 전부는 기후 격자 셀 사이의 편향 보정에서 왔다. 이 연구의 기여는 평가 설계, 음성 결과를 포함한 그 결과, 그리고 그 결과로 만든 라벨 수별 워크플로다.
